% LaTeX Curriculum Vitae Template
%
% Copyright (C) 2004-2009 Jason Blevins <jrblevin@sdf.lonestar.org>
% http://jblevins.org/projects/cv-template/
%
% You may use use this document as a template to create your own CV
% and you may redistribute the source code freely. No attribution is
% required in any resulting documents. I do ask that you please leave
% this notice and the above URL in the source code if you choose to
% redistribute this file.

\documentclass[letterpaper]{article}

\usepackage{hyperref}
\usepackage{geometry}
\usepackage{xcolor}

% Comment the following lines to use the default Computer Modern font
% instead of the Palatino font provided by the mathpazo package.
% Remove the 'osf' bit if you don't like the old style figures.
\usepackage[T1]{fontenc}
\usepackage[sc,osf]{mathpazo}

% Set your name here
\def\name{Nesar Ramachandra}

% Replace this with a link to your CV if you like, or set it empty
% (as in \def\footerlink{}) to remove the link in the footer:
%\def\footerlink{http://jblevins.org/projects/cv-template/}

% The following metadata will show up in the PDF properties
\hypersetup{
  colorlinks = true,
  urlcolor = black,
  pdfauthor = {\name},
  pdfkeywords = {},
  pdftitle = {\name: Curriculum Vitae},
  pdfsubject = {Curriculum Vitae},
  pdfpagemode = UseNone
}

\geometry{
  body={6.5in, 8.5in},
  left=1.0in,
  top=1.25in
}

% Customize page headers
\pagestyle{myheadings}
\markright{\name}
\thispagestyle{empty}

% Custom section fonts
\usepackage{sectsty}
\sectionfont{\rmfamily\mdseries\Large}
\subsectionfont{\rmfamily\mdseries\itshape\large}

% Other possible font commands include:
% \ttfamily for teletype,
% \sffamily for sans serif,
% \bfseries for bold,
% \scshape for small caps,
% \normalsize, \large, \Large, \LARGE sizes.

% Don't indent paragraphs.
\setlength\parindent{0em}

% Make lists without bullets
\renewenvironment{itemize}{
  \begin{list}{}{
    \setlength{\leftmargin}{1.5em}
  }
}{
  \end{list}
}

\begin{document}

% Place name at left
{\huge \name}

% Alternatively, print name centered and bold:
%\centerline{\huge \bf \name}

\vspace{0.25in}

\begin{minipage}{0.45\linewidth}
  \href{https://www.anl.gov/cps}{Computational Science Division} \\
  \href{https://www.anl.gov}{Argonne National Laboratory} \\
  9700 Cass Avenue, Lemont, IL 60439
\end{minipage}
\begin{minipage}{0.45\linewidth}
  \begin{tabular}{ll}
    % Phone: & +1 785-979-6181 \\
    Email: & \href{mailto:nramachandra@anl.gov}{\tt nramachandra@anl.gov} \\
    Homepage: & \href{https://nesar.github.io/}{\tt nesar.github.io} \\
  \end{tabular}
\end{minipage}


% \section*{Current position}
\section*{Professional experience}

\begin{itemize}
\item Assistant Computational Scientist, Argonne National Laboratory, \textbf{Aug, 2021 - Present.}
\item Adjunct Faculty Member, University of Houston, \textbf{Aug, 2025 - Present.}
\item Postdoctoral Researcher -  Cosmological Physics and Advanced Computing Group, High Energy Physics Division, Argonne National Laboratory, \textbf{Aug, 2018 - Aug, 2021.}
\item Associate Fellow - Kavli Institute for Cosmological Physics, University of Chicago, \textbf{March, 2019 - January, 2024.}
\item Fellow - Kavli Summer Program on Astrophysics, University of California, Santa Cruz, \textbf{June, 2019 - Aug, 2019.}
\item Graduate researcher - Cosmological Physics and Advanced Computing (CPAC) Group, Argonne National Laboratory, \textbf{Jan, 2018 - Aug, 2018.}
\item HEP-Center for Computational Excellence Fellow - High Energy Physics Division, Argonne National Laboratory, \textbf{June, 2017 - Aug, 2017.}
\item Head Teaching Assistant, Department of Physics and Astronomy, University of Kansas, \textbf{Aug, 2015 - May, 2016.}
\item Graduate Research Assistant -  Department of Physics and Astronomy, University of Kansas, \textbf{Aug, 2013 - June, 2017.}
% \item Graduate Teaching Assistant, Department of Physics and Astronomy, University of Kansas, Aug, 2014 - May, 2015
\item Research Scholar - Tata Institute of Fundamental Research Centre for Interdisciplinary Science, Hyderabad, India, \textbf{Jan, 2013 - July, 2013.}
\item Research scholar - Indian Institute of Astrophysics, Bangalore, India, \textbf{Jan 2012 - Jan 2013.}
\end{itemize}

\section*{Academic background}

\begin{enumerate}
  \item Ph.D. (Hons.) Physics, University of Kansas, \textbf{2018.}
\begin{itemize} 
  \item Dissertation title: {\it Topology, Geometry and Morphology of the Dark Matter web.}
\item Adviser: Professor Sergei Shandarin
\end{itemize}
  \item Integrated M.S. (Hons.) Physics, Birla Institute of Technology and Science (BITS) Pilani, India, \textbf{2012.}
\begin{itemize}   
  \item Thesis title: {\it Dynamics of ellipsoidal collapse in a cosmological setting.}
    \item Adviser: Professor Arun Mangalam, Indian Institute of Astrophysics.
\end{itemize}
\end{enumerate}

\section*{Research interests and experience}

\textbf{Techniques:} Scientific machine learning and foundation models, high-performance computing, numerical methods, Bayesian inference, topological and geometrical techniques, agentic AI systems. 

\vspace{0.05in}
\textbf{Domains:} Astrophysics, Nuclear and Particle physics, Data analysis experience in Climate science, and aero-structure manufacturing. 

% \section*{Projects}


\section*{Publications}

% Citation record \href{https://scholar.google.co.in/citations?user=dCe3MK0AAAAJ&hl=en}{(Google Scholar profile)}: \qquad  Total citations: 838 \qquad h-Index: 11 

% \subsection*{In preparation}

% \begin{enumerate}

% \item \textbf{N. Ramachandra}, A. Fadikar, J. Chaves-Montero, S. Habib, K. Heitmann: Photometric Redshift Inference using Synthetic Training of Gaussian Mixture Models.

% \item K. Fukami, R. Maulik, \textbf{N. Ramachandra}, K. Taira, K. Fukagata: Voronoi CNNs: Interfacing unstructured data with convolutional architectures for fluid flow super resolution.

% \item A. Hearin, \textbf{N. Ramachandra}, J. DeRose, M. Becker: SHAMnet: Differentiable Abundance Matching.

% \item \textbf{N. Ramachandra}, M. Binois, S. Habib, K. Heitmann: Gaussian Process Emulation of CMB Power Spectra with Autoencoder-based Dimensional Reduction.

% \end{enumerate}


\subsection*{Under review}

\begin{enumerate}

\item LSST Dark Energy Science Collaboration {\it et al.} (incl. \textbf{N. Ramachandra}): Opportunities in AI/ML for the Rubin LSST Dark Energy Science Collaboration, Submitted, {\it arxiv:2601.14235}, 2026.

\item J. Li, Y. Ting, A. Accomazzi, T. Ghosal, \textbf{N. Ramachandra}: Predicting New Concept-Object Associations in Astronomy by Mining the Literature, Submitted, {\it arxiv:2602.14335}, 2026.

\item \textbf{N. Ramachandra}, N. Frontiere, M. Buehlmann, K. R. Moran, J. D. Emberson, K. Heitmann, S. Habib: Emulator-Based Inference of Cosmological Subgrid Models, Submitted, {\it arxiv:2601.07306}, 2026.

\item Y. Comlek {\it et al.} (incl. \textbf{N. Ramachandra}): Multi-task Modeling for Engineering Applications with Sparse Data, Submitted, {\it arxiv:2601.05910}, 2026.

\item N. Frontiere, J. D. Emberson, M. Buehlmann, S. Habib, K. Heitmann, \textbf{N. Ramachandra}, C. Faucher-Gigu\`ere: Modeling Galaxy Formation in Cosmological Simulations with CRK-HACC, Submitted, {\it arxiv:2511.21921}, 2025.

\item J. J. Bock {\it et al.} (incl. \textbf{N. Ramachandra}): The SPHEREx Satellite Mission, Submitted, {\it arxiv:2511.02985}, 2025.

\item X. Dong, \textbf{N. Ramachandra}, S. Habib, K. Heitmann: Benchmarking AI-evolved cosmological structure formation, Submitted, {\it arxiv:2510.06731}, 2025.

\item T. de Haan {\it et al.} (incl. \textbf{N. Ramachandra}): AstroMLab 4: Benchmark-Topping Performance in Astronomy Q\&A with a 70B-Parameter Domain-Specialized Reasoning Model, Submitted, {\it arxiv:2505.17592}, 2025.

\item F. Cappello {\it et al.} (incl. \textbf{N. Ramachandra}): EAIRA: Establishing a Methodology for Evaluating AI Models as Scientific Research Assistants, Submitted to the International Journal of High Performance Computing Applications {\it arxiv:2502.20309}, 2025. 


\item M. Michael Rau, F. Kéruzoré, \textbf{N. Ramachandra}, L. Bleem: Reducing Model Error Using Optimised Galaxy Selection: Weak Lensing Cluster Mass Estimation, Submitted to the Monthly Notices of the Royal Astronomical Society, {\it arxiv:2406.11950}.


% \item D. Huppenkothen, M. Ntampaka, M. Ho, M. Fouesneau, B. Nord, J. E. G. Peek, M. Walmsley, J. F. Wu, C. Avestruz, T. Buck, M. Brescia, D. P. Finkbeiner, A. D. Goulding, T. Kacprzak, P. Melchior, M. Pasquato, \textbf{N. Ramachandra}, Y.-S. Ting, G. van de Ven, S. Villar, V.A. Villar, E. Zinger: Constructing Impactful Machine Learning Research for Astronomy: Best Practices for Researchers and Reviewers, Submitted, Bulletin of the American Astronomical Society, {\it arxiv:2310.12528}

\item D. Huppenkothen {\it et al.} (incl. \textbf{N. Ramachandra}): Constructing Impactful Machine Learning Research for Astronomy: Best Practices for Researchers and Reviewers, Submitted, Bulletin of the American Astronomical Society, {\it arxiv:2310.12528}
 




\item T. Y. Chen, B. Dey, A. Ghosh, M. Kagan, B. Nord, \textbf{N. Ramachandra}: Interpretable Uncertainty Quantification in AI for HEP, Submitted to the Proceedings of the US Community Study on the Future of Particle Physics (Snowmass 2021), {\it arxiv:2208.03284}. 


\item S. Madireddy, \textbf{N. Ramachandra}, N. Li, , P. Balaprakash, S. Habib: A Modular Deep Learning Pipeline for Galaxy-Scale Strong Gravitational Lens Detection and Modeling, Submitted to the Open Journal of Astrophysics, {\it arxiv:1911.03867v3}.



\end{enumerate}

\subsection*{Refereed journal articles}
% (peer-reviewed)

\begin{enumerate}


\item A. Ganguli, \textbf{N. Ramachandra}, N., Bessac, J., Constantinescu, E.: Enhancing interpretability in generative modeling: statistically disentangled latent spaces guided by generative factors in scientific datasets. Machine Learning, 114(9), 197., {\it https://arxiv.org/abs/2507.00298}


\item T. de Haan {\it et al.} (incl. \textbf{N. Ramachandra}): AstroMLab 3: Achieving GPT-4o Level Performance in Astronomy with a Specialized 8B-Parameter Large Language Model, Scientific Reports, Volume 15, Issue 1, id.13751, {\it arxiv:2411.09012}


\item Anirban Samaddar, S. K. Ravi, \textbf{N. Ramachandra}, L. Luan,  S. Madireddy, A. Bhaduri, P. Pandita, C. Sun, L. Wang: Data-Efficient Dimensionality Reduction and Surrogate Modeling of High-Dimensional Stress Fields, Accepted at The ASME Journal of Mechanical Design.



\item Y. Ting {\it et al.} (incl. \textbf{N. Ramachandra}): AstroMLab 1: Who Wins Astronomy Jeopardy!?, Astronomy and Computing, Volume 51, id.100893, {\it arxiv:2407.11194}

\item A. Bhaduri, \textbf{N. Ramachandra}, S. K. Ravi, L. Luan, P. Pandita, P. Balaprakash, M. Anitescu, C. Sun, L. Wang: Efficient mapping between void shapes and stress fields using Deep Convolutional Neural Networks with sparse data, ASME Selected Papers from CIE2023, ASME Journal of Computing and Information Science in Engineering (JCISE), 2024 

\item H. Wang, S. Sreejith, Y. Lin, \textbf{N. Ramachandra}, A. Slosar, S. Yoo: Neural Network Based Point Spread Function Deconvolution For Astronomical Applications, {\it Open Journal of Astrophysics}, 2023.

\item M. Lucey, N. A. Kharusi, K. Hawkins, Y. Ting, \textbf{N. Ramachandra}, A. M. Price-Whelan, T. C. Beers, Y. S. Lee, J. Yoon: Carbon-enhanced metal-poor star candidates from BP/RP spectra in Gaia DR3, {\it Monthly Notices of the Royal Astronomical Society}, 2023.


\item \textbf{N. Ramachandra}, J. Chaves-Montero, A. Alarcon, A. Fadikar, S. Habib, K. Heitmann: Machine learning synthetic spectra for probabilistic redshift estimation: SYTH-Z, {\it Monthly Notices of the Royal Astronomical Society}, 2022.

\item A. Hearin, \textbf{N. Ramachandra}, M. Becker, J. DeRose: Differentiable Predictions for Large Scale Structure with SHAMNet, {\it Open Journal of Astrophysics}, 2021.

\item K. Fukami, R. Maulik, \textbf{N. Ramachandra}, K. Fukagata, K. Taira : Global field reconstruction from sparse sensors with Voronoi tessellation-assisted deep learning, {\it Nature Machine Intelligence}, 2021.


\item L. Vazsonyi, P. Taylor, G. Valogiannis, \textbf{N. Ramachandra}, A. Ferte, J. Rhodes: Constraining f(R) Gravity with a k-cut Cosmic Shear Analysis of the Hyper Suprime-Cam First-Year Data, {\it Physical Review D}, 2021. 

\item K. Storey-Fisher, M. Huertas-Company, \textbf{N. Ramachandra}, F. Lanusse, A. Leauthaud, Y. Luo, S. Huang,  X. Prochaska: Anomaly detection in Hyper Suprime-Cam galaxy images with generative
adversarial networks, {\it Monthly Notices of the Royal Astronomical Society}, 2021.

\item Y. Wang, \textbf{N. Ramachandra}, E. Salazar-Canizales, H. Feldman, R. Watkins, K. Dolag: Peculiar Velocity Estimation from Kinetic SZ Effect using Deep Neural Networks, Accepted, {\it Monthly Notices of the Royal Astronomical Society}, 2021.

\item \textbf{N. Ramachandra}, G. Valogiannis, M. Ishak, K. Heitmann (for the LSST Dark Energy Science Collaboration): Matter Power Spectrum Emulator for f(R) Modified Gravity Cosmologies, Accepted, {\it Physical Review D}, 2021.
% {\it arXiv:2010.00596}. 

\item T. Cheng, M. Huertas-Company, C. Conselice, A. Aragón-Salamanca, B. Robertson, \textbf{N. Ramachandra}: Beyond the Hubble Sequence -- Exploring Galaxy Morphology with Unsupervised Machine Learning, {\it Monthly Notices of the Royal Astronomical Society}, Volume 503, Issue 3, May 2021, Pages 4446–4465.

\item R. Maulik, T. Botsas, \textbf{N. Ramachandra}, M. Lachlan, I. Pan: Latent-space time evolution of non-intrusive reduced-order models using Gaussian process emulation, Accepted, {\it Physica D: Nonlinear Phenomena}, 2020.





\item R. Maulik, K. Fukami, \textbf{N. Ramachandra}, K. Fukagata, K. Taira : Probabilistic neural networks for fluid flow model-order reduction and data recovery, {\it Physical Review Fluids}, 5, 104401, 2020.

\item M. Lucey, Y. Ting, \textbf{N. Ramachandra}, K. Hawkins: From the Inner to Outer Milky Way: A Photometric Sample of 2.6 Million Red Clump Stars, {\it Monthly Notices of the Royal Astronomical Society}, Volume 495, Issue 3, Pages 3087–3103.

\item L. Bleem {\it et al.} (incl. \textbf{N. Ramachandra}): The SPTpol Extended Cluster Survey, {\it The Astrophysical Journal Supplement Series}, Volume 247, Number 1.

% \item L. Bleem, S. Bocquet, B. Stalder, M. Gladders, P. Ade, S. Allen, A. Anderson, J. Annis, M. Ashby, J. Austermann, S. Avila, J. Avva, M. Bayliss, J. Beall, K. Bechtol, A. Bender, B. Benson, E. Bertin, F. Bianchini, C. Blake, M. Brodwin, D. Brooks, E. Buckley-Geer, D. Burke, J. Carlstrom, A. Rosell, M. Kind, J. Carretero, C. Chang, H. Chiang, R. Citron, C. Moran, M. Costanzi, T. Crawford, A. Crites, L. Costa, T. Haan, J. Vicente, S. Desai, H. Diehl, J. Dietrich, M. Dobbs, T. Eifler, W. Everett, B. Flaugher, B. Floyd, J. Frieman, J. Gallicchio, J. García-Bellido, E. George, D. Gerdes, A. Gilbert, D. Gruen, R. Gruendl, J. Gschwend, N. Gupta, G. Gutierrez, N. Halverson, N. Harrington, J. Henning, C. Heymans, G. Holder, D. Hollowood, W. Holzapfel, K. Honscheid, J. Hrubes, N. Huang, J. Hubmayr, K. Irwin, D. James, T. Jeltema, S. Joudaki, G. Khullar, M. Klein, L. Knox, N. Kuropatkin, A. Lee, D. Li, C. Lidman, A. Lowitz, N. MacCrann, G. Mahler, M. Maia, J. Marshall, M. McDonald, J. McMahon, P. Melchior, F. Menanteau, S. Meyer, R. Miquel, L. Mocanu, J. Mohr, J. Montgomery, A. Nadolski, T. Natoli, J. Nibarger, G. Noble, V. Novosad, S. Padin, A. Palmese, D. Parkinson, S. Patil, F. Paz-Chinchón, A. Plazas, C. Pryke, \textbf{N. Ramachandra}, C. Reichardt, J. Remolina González, A. Romer, A. Roodman, J. Ruhl, E. Rykoff, B. Saliwanchik, E. Sanchez, A. Saro, J. Sayre, K. Schaffer, T. Schrabback, S. Serrano, K. Sharon, C. Sievers, G. Smecher, M. Smith, M. Soares-Santos, A. Stark, K. Story, E. Suchyta, G. Tarle, C. Tucker, K. Vanderlinde, T. Veach, J. Vieira, G. Wang, J. Weller, N. Whitehorn, W. Wu, V. Yefremenko, Y. Zhang: The SPTpol Extended Cluster Survey, {\it The Astrophysical Journal Supplement Series}, Volume 247, Number 1. 


\item \textbf{N. Ramachandra}: Topology, Geometry and Morphology of the Dark Matter Web (Doctoral Dissertation), {\it ProQuest Dissertations \& Theses Global database}, Publication Number: 10845180. 

\item N. Libeskind {\it et al.} (incl. \textbf{N. Ramachandra}): Tracing the cosmic web, {\it Monthly Notices of the Royal Astronomical Society}, Volume 473, Issue 1, Pages 1195-1217.


% \item N. Libeskind, R. Weygaert, M. Cautun, B. Falck, E. Tempel, T. Abel, M. Alpaslan, M. Aragón-Calvo, J. Forero-Romero, R. Gonzalez, S.  Gottlöber, O. Hahn, W. Hellwing, Y. Hoffman, B. Jones, F. Kitaura, A. Knebe, S. Manti, M. Neyrinck, S. Nuza, N. Padilla, E. Platen, \textbf{N. Ramachandra}, A. Robotham, E. Saar, S. Shandarin, M. Steinmetz, R. Stoica, T. Sousbie, G. Yepes: Tracing the cosmic web, {\it Monthly Notices of the Royal Astronomical Society}, Volume 473, Issue 1, Pages 1195-1217.

\item \textbf{N. Ramachandra}, S. Shandarin: Dark matter haloes: a multistream view, {\it Monthly Notices of the Royal Astronomical Society}, Volume 470, Issue 3, p. 3359-3373.

\item \textbf{N. Ramachandra}, S. Shandarin: Topology and geometry of the dark matter web: a multistream view, {\it Monthly Notices of the Royal Astronomical Society}, Volume 467, Issue 2, p.1748-1762.

\item \textbf{N. Ramachandra}, S. Shandarin: Multi-stream portrait of the Cosmic web, {\it Monthly Notices of the Royal Astronomical Society}, Volume 452, Issue 2, p.1643-1653. 

\end{enumerate}



\subsection*{Papers in technical conference proceedings – refereed}

\begin{enumerate}


\item T. Alghamdi, J. Xu, \textbf{N. Ramachandra}, N. Sato, Y. Li: Towards an event-level analysis in hadronic physics using generative AI-based surrogates, Accepted for the proceedings of the IEEE International
Conference on Tools with Artificial Intelligence, 2025. 

\item B. Xia, \textbf{N. Ramachandra}, A. Wells, S. Habib, J. Wise: Multi-modal Foundation Model for Cosmological
Simulation Data, Accepted at the NeurIPS ML4Physics workshop, 2025. 


\item A. Ganguli, A. Samaddar, F. K\'eruzor\'e \textbf{N. Ramachandra}, J. Bessac, S. Madireddy, E. Constantinescu: Uncovering Physical Drivers of Dark Matter Halo Structures with Auxiliary- Variable-Guided Generative Models, Accepted at NeurIPS ML4Physics workshop, 2025. 

\item J. Z. Tam {\it et al.} (incl. \textbf{N. Ramachandra}): InferA: A Smart Assistant for Cosmological Ensemble Data, Accepted at the International Conference for High Performance Computing, Networking, Storage, and Analysis (SC25). 


\item \textbf{N. Ramachandra}, Ting, Y. S., Sun, Z., Wells, A., Habib, S.: Teaching LLMs to Speak Spectroscopy, Accepted at the ICML ML4Astro workshop, 2025. 


\item T. de Haan {\it et al.} (incl. \textbf{N. Ramachandra}): AstroSage: Leading Performance in Astronomy Q\&A with a 70B-Parameter Domain-Specialized
Model, Accepted at the ICML ML4Astro workshop, 2025.


\item A. Ganguli, \textbf{N. Ramachandra}, J. Bessac, E. Constantinescu: Disentangling Latent Spaces in Generative Models for Scientific Datasets, Accepted at Science Understanding through Data Science Conference-2024.


\item A. Wells, \textbf{N. Ramachandra}, S. Habib,  N. Frontiere: Learning Relationships Between Disparate Representations of Objects with Transformers and Contrastive Losses, Accepted at Science Understanding through Data Science Conference-2024.

\item L. Luan, S. Krishnan Ravi, A. Bhaduri, P. Pandita, L. Wang, \textbf{N. Ramachandra}, S. Madireddy: High-dimensional Surrogate Modeling for Image Data with Nonlinear Dimension Reduction, Accepted at {\it American Institute of Aeronautics and Astronautics SciTech Forum 2024}
% https://arc.aiaa.org/doi/10.2514/6.2024-0388

\item L. Luan, \textbf{N. Ramachandra}, S. K. Ravi, A. Bhaduri, P. Pandita, P. Balaprakash, M. Anitescu, C. Sun, L. Wang: Application of probabilistic modeling and automated machine learning framework for high-dimensional stress field, Accepted at {\it ASME-IDETC Conference 2023}, {\it arxiv:2303.16869}.


\item A. Bhaduri, \textbf{N. Ramachandra}, S. K. Ravi, L. Luan, P. Pandita, P. Balaprakash, M. Anitescu, C. Sun, L. Wang: Efficient mapping between void shapes and stress fields using Deep Convolutional Neural Networks with sparse data, Accepted at {\it ASME-IDETC Conference 2023.}

\item H. Wang, S. Sreejith, Y. Lin, \textbf{N. Ramachandra}, A. Slosar, S. Yoo: Deconvolution of Astronomical Images with Deep Neural Networks, Accepted at {\it NeurIPS workshop on AI for Science}, 2022.

\item X. Dong, \textbf{N. Ramachandra}, S. Habib, K. Heitmann, M. Buehlmann, S. Madireddy: Physical Benchmarking for AI-Generated Cosmic Web, Accepted at {\it NeurIPS workshop on AI for Science}, 2021.

\item K. Storey-Fisher, M. Huertas-Company, \textbf{N. Ramachandra}, F. Lanusse, A. Leauthaud, Y. Luo, S. Huang: Anomaly Detection in Astronomical Images with Generative Adversarial Networks, Accepted at {\it NeurIPS workshop on Machine Learning and the Physical Sciences}, 2020.

\item K. Fukami, R. Maulik, \textbf{N. Ramachandra}, K. Fukagata, K. Taira: Probabilistic neural network-based reduced-order surrogate for fluid flows, Accepted at {\it NeurIPS workshop on Machine Learning and the Physical Sciences}, 2020.

\item S. Madireddy, N. Li, \textbf{N. Ramachandra}, P. Balaprakash, S. Habib: Modular Deep Learning Analysis of Galaxy-Scale Strong Lensing Images, {\it Machine Learning and the Physical Sciences Workshop at the 33rd Conference on Neural Information Processing Systems (2019)}.

\end{enumerate}

% \subsection*{Papers in technical conference proceedings – non-refereed}


% \subsection*{Technical reports}


\section*{External software developed}


\begin{enumerate}

\item \href{https://github.com/LSSTDESC/mgemu}{\tt MGemu: The Modified Gravity Emulator}, \textbf{N. Ramachandra}, G. Valogiannis, M. Ishak, K. Heitmann (for the LSST Dark Energy Science Collaboration).

\item \href{https://github.com/nesar/MDN_phoZ}{\tt SYTH-Z: The galactic redshift estimator}, \textbf{N. Ramachandra}, J. Chaves-Montero, A. Alarcon, A. Fadikar, S. Habib, K. Heitmann.

\end{enumerate}

\section*{Talks}
\subsection*{Major Conferences and Symposia}
% (Invited talks at  presented by the candidate)}

\begin{enumerate}

\item \textbf{N. Ramachandra}: Emulator-Based Inference of Cosmological Subgrid Models in HACC, Mock NYC 2026: Improving Galaxy Modeling for Cosmological Inference, Flatiron Institute, New York City, 2026.

\item \textbf{N. Ramachandra}, J. Coburn, A. Wells : Multi-agent systems for Physics research, JLESC: AI for Science, Chicago, 2025. 

\item \textbf{N. Ramachandra} : Astrophysics benchmarks for language and reasoning models, TPC25 San Jose, 2025. 

\item \textbf{N. Ramachandra}, A. Ganguli, J. Bessac, E. Constantinescu: Latent Space Modeling for AI-Enabled Physics Research, SIAM CSE - Generative Models for Scientific Applications Dallas, 2025.

\item \textbf{N. Ramachandra}: Foundation models and Agents for Physics, CFNS workshop on Probing the frontiers of nuclear physics with AI at the EIC, Stony Brook University, NY, 2025. 

\item \textbf{N. Ramachandra}, N. Frontiere, K. Moran, K. Heitmann, S. Habib : Subgrid Model Calibration for Cosmology Simulations, Cosmology and galaxy astrophysics with simulations and machine learning, Flatiron Institute, New York City, 2024 

\item \textbf{N. Ramachandra}, N. Frontiere, K. Moran, K. Heitmann, S. Habib : Sub-grid parameter calibration for Cosmological
simulations, 2024 SciDAC Principal Investigator (PI) Meeting, Washington DC, 2024.

\item \textbf{N. Ramachandra}, A. Ganguli, J. Bessac, E. Constantinescu, Foundation models for physical sciences, QCD at the Femtoscale in the Era of Big Data, University of Washington, 2024.

\item \textbf{N. Ramachandra}, A. Ganguli, J. Bessac, E. Constantinescu, Latent spaces in generative models, QCD at the Femtoscale in the Era of Big Data, University of Washington, 2024.

\item \textbf{N. Ramachandra}, Physics Benchmarking strategies for Large Language Models, AuroraGPT Evaluation strategy meeting, 2024.

\item \textbf{N. Ramachandra}, M. Buehlmann, S. Habib, P. Larsen: Cosmology and EBL simulations, SPHEREx Science L4 SIR Peer Review and Science Team Meeting, Pasadena, 2023.

\item \textbf{N. Ramachandra}, K. Moran, N. Frontiere, K. Heitmann, S. Habib : Calibrating sub-grid hydrodynamical parameters
using cosmic emulators, 2023 SciDAC Principal Investigator (PI) Meeting, Washington DC, 2023.


\item \textbf{N. Ramachandra}, H. Wang, S. Sreejith, Y. Lin, A. Slosar, S. Yoo: Deconvolution of Astronomical Images with Deep Neural Networks, NeurIPS, 2023.

\item \textbf{N. Ramachandra}: Estimation of Galaxy Redshift with Probabilistic Neural Networks, SIAM Conference on Parallel Processing for Scientific Computing, Seattle, 2022 (remote).


\item \textbf{N. Ramachandra}: Astrophysical AI: Issues and lessons learned, Exascale Astronomy AI and benchmarking workshop, Edinburgh, 2020 (remote).  

\item \textbf{N. Ramachandra}, G. Valogiannis, M. Ishak, K. Heitmann: Matter Power Spectrum Emulator for f(R) Modified Gravity Cosmologies, Science Highlight Plenary talk, LSST-DESC July Meeting, 2020. 

\item \textbf{N. Ramachandra}: Scientific Machine Learning using Astrophysical Simulations, Conference on Data Analysis, Santa Fe, 2020.

\item \textbf{N. Ramachandra}: Scientific machine learning with synthetic astrophysical data, Machine Learning Tools for Research in Astronomy Ringberg, Germany, 2019.

\item \textbf{N. Ramachandra}, S. Habib, K. Heitmann: Cosmological analysis pipelines through Neural Networks, American Physical Society April Meeting, Columbus, Ohio 2018.

\item \textbf{N. Ramachandra}, S. Shandarin: Topology and geometry of the dark matter web, American Physical Society April Meeting, Washington D.C, 2017

\item \textbf{N. Ramachandra}, S. Shandarin: The Multi-stream portrait of the cosmic web, American Physical Society April Meeting, Salt Lake City, Utah, 2016. 

\item \textbf{N. Ramachandra}, S. Shandarin: The Multi-stream structures of the cosmic web, Canadian-American-Mexican Graduate Students Physics Conference at Oaxaca, Mexico, 2015.

\end{enumerate}

\subsection*{Seminars and Colloquia}
% presented by the candidate}

%(especially external to Argonne)

\begin{enumerate}

\item \textbf{N. Ramachandra}, K. Heitmann: The Knowledge Extraction IDA HEP Pilot Project, LSST-DESC Seminar, 2026.

\item \textbf{N. Ramachandra}, S. Habib: Synthetic sky catalogs, NASA-SPHEREx meeting, California Institute of Technology, Pasadena, 2023.


\item \textbf{N. Ramachandra}, P. Balaprakash: DeepHyper at scale: NERSC Perlmutter, ALCF DeepHyper Automated Machine Learning Workshop, Lemont, 2022


\item \textbf{N. Ramachandra}: Cosmology using Scientific Machine Learning, AI \& HPC Seminar, Argonne National Laboratory, 2020.

\item \textbf{N. Ramachandra}, S. Madireddy, N. Li, P. Balaprakash, S. Habib: Probabilistic Deep Learning for Galaxy-Scale Strong Lensing Studies, University of Chicago, 2020.



% CPS seminar
\item \textbf{N. Ramachandra}: Deep Generative models and Astrophysical Image emulation, Machine Learning in the Era of Large Astronomical Surveys, Santa Cruz, 2019.

\item \textbf{N. Ramachandra}, G. Valogiannis, M. Ishak, K. Heitmann: Gaussian Process Emulators for Modified Gravity Summary Statistics, LSST-DESC meeting, Paris, France, 2019 (remote).

\item \textbf{N. Ramachandra}, M. Binois, S. Habib, K. Heitmann: Suite of Gaussian Process emulators for cosmological inference problems, Likelihood-free inference workshop, Flatiron Institute, New York, 2019.

\item \textbf{N. Ramachandra}, M. Binois, S. Habib, K. Heitmann: Cosmology meets Machine learning: Emulation of CMB Power Spectra, Astro seminar, University of Kansas, Lawrence, 2019.

\item \textbf{N. Ramachandra}, M. Binois, S. Habib, K. Heitmann: Variational autoencoders for the emulation of cosmological functions, KICP Postdoc Symposium, University of Chicago, 2019. 

\item \textbf{N. Ramachandra}, P. Larsen: Cosmic Emulators for next generation surveys, Accurate lensing in the era of precision Cosmology, University of California, Berkeley, 2019. 

\item \textbf{N. Ramachandra}, S. Madireddy, N. Li, P. Balaprakash, S. Habib: Deep learning pipelines for lensing analysis, Astrophysics Seminar, University of Kansas, 2017.

\item \textbf{N. Ramachandra}, S. Madireddy, N. Li, P. Balaprakash, S. Habib: Strong Lensing analysis using Deep Neural Networks, Young Scientists Symposium, Argonne National Laboratory, 2017. 

\item \textbf{N. Ramachandra}: Emulation of the halo mass function, SAMSI Research Triangle Park, North Carolina, 2017

\item \textbf{N. Ramachandra}, S. Shandarin: The dynamical structure of the cosmic web, MidAmerican Regional Astrophysics Conference at The University of Missouri, 2015.

\end{enumerate}


\subsection*{General-interest lectures}
% \subsection*{Other presentations (contributed talks, general-interest lectures, etc.)}

\begin{enumerate}

\item \textbf{N. Ramachandra}: Artificial Intelligence assisted understanding of the Universe, Art of Science Public lecture, Chicago, 2023.

\item \textbf{N. Ramachandra}: Cosmology in the era of Artificial Intelligence, Lifelong Learning Outreach Series, Chicago, 2021. 

\end{enumerate}


% \subsection*{Chaired/organization of national/international workshops/technical meetings/conferences }

\section*{Organization}

\begin{enumerate}

\item Organizer: Minisymposium on Generative Models for Scientific Applications, SIAM CSE, 2025.

\item Co-organizer: ALCF DeepHyper Automated Machine Learning Workshop, 2022

\item Co-organizer: AI, Statistics and Machine Learning Journal Club, Argonne National Laboratory, 2019. 

\item Co-organizer: Workshop on Advanced Statistical Methods Meet Machine Learning, Argonne National Laboratory, 2018-2019 

\end{enumerate}



% \subsection*{Developed/managed/contributed to projects}
\section*{Major projects}

\begin{enumerate}

\item Member, AuroraGPT -- A large-scale AI models for Science, Argonne National Laboratory, 2024.

\item Principal Investigator, Laboratory Computing Resource Center (LCRC) -- AI-modeling of cosmic fields and dark matter halos, Argonne National Laboratory, 2020-2024.

\item Principal Investigator, Laboratory Directed Research and Development (LDRD) -- Pushing the mapping limits of the cosmological evolution, Argonne National Laboratory, 2023. 

\item Co-Principal Investigator, SciDAC-5 -- Femtoscale Imaging of Nuclei using Exascale Platforms, DOE-NP, 2022.

\item Co-Principal Investigator, Probabilistic Machine Learning for Rapid Large-Scale and High-Rate Aerostructure Manufacturing, with Argonne, GE Research, Edison Welding Institute and GKN Aerospace, DOE-EERE, 2021.  
 
\item Member: SciDAC-5 -- Enabling Cosmic Discoveries in the Exascale Era, DOE-HEP, 2022.

\item Principal Investigator, Laboratory Directed Research and Development (LDRD) -- Artificial Intelligence (AI)-enabled analytics for Cosmology, Argonne National Laboratory, 2022-2024.  

\item Principal Investigator, Argonne Laboratory Computing Resource Center allocation -- AI-modeling of cosmic fields and dark matter halos, 2022-2023.
 
\item Co-Principal Investigator, Swift LDRD: Deep-Learning Inference for Nuclear Femtography on Exascale Platforms, Argonne National Laboratory, 2022. 

\item Member, NASA-SPHEREx mission L4 Science Team, 2020-2024 (launched March 2025)
 
\item Co-Principal Investigator, Laboratory Directed Research and Development (LDRD) Expedition Project: Deep Learning-Based Scalable and Robust Strong Gravitational Lensing Characterization Pipeline Using SambaNova, Argonne National Laboratory, 2021.
  
\item Co-Principal Investigator, Laboratory Directed Research and Development (LDRD): Automated Model Inference for Cosmological Structure Formation using Reinforcement Learning, Argonne National Laboratory, 2021. 
  
\item Co-lead, Strong Lensing Working group project, LSST-DESC: Deep Learning to Identify and Deblend Strong Lenses for LSST. 

\item Co-lead, Photometric Redshifts Working group project, LSST-DESC: Synthetic-spectra for sample redshift inference  
  
\item Co-lead: Modeling and Combined Probes (MCP) Working group project, LSST-DESC, Power spectra emulators for $f(R)$ Modified Gravity Cosmologies.
  
\item Member: SciDAC-4 -- Inference and Machine Learning at Extreme Scales, DOE-HEP.


\end{enumerate}


\section*{Professional service}

\begin{enumerate}

% \item Journal referee, Physica D: Nonlinear Phenomena, Elsevier Publishing, 2023. 

\item Journal referee, Astronomy and Computing, Science Direct, 2025

\item Journal referee, Monthly Notices of Royal Astronomical Society - Oxford University Press, 2018, 2023, 2025. 

\item Journal referee, JMIR Medical Informatics, 2025. 

\item International Journal of Intelligent Systems -  The Wiley Hindawi publishing, 2023.  

\item Conference proceedings referee, NeurIPS workshop on Machine Learning and the Physical Sciences, 2021, 2023, 2025. 

\item Journal referee, The Astrophysical Journal - IOPscience, 2021. %2020

\item Journal referee, Astronomy and Computing - Elsevier Publishing, 2019.

\item Journal Referee, Journal of Cosmology and Astroparticle Physics - IOP Publishing, 2015. 

\item Grant reviewer: Future Investigators in NASA Earth and Space Science and Technology (FINESST) program, 2019-2020.

\end{enumerate}



\section*{Teaching experience}
\begin{enumerate}

\item {\it Head Teaching Assistant}, Department of Physics and Astronomy, University of Kansas (2015-2016): Implemented, oversaw and evaluated new teaching strategies for 12 undergraduate courses.
\item {\it Graduate Teaching Assistant}, Department of Physics and Astronomy, University of Kansas (2014-2015): General Physics II Honors, General Physics II Honors, General Physics I, College Physics I.
\item {\it Undergraduate Teaching Assistant}, Department of Physics, BITS-Pilani (2011): Theory of Relativity.
 
 
\end{enumerate}

% \section*{Mentorship}

% \begin{enumerate}

%   \item Postdoc researchers: Arkaprabha Ganguli (Argonne), Anirban Samaddar (Argonne)

%   \item Summer internship students: Claire Guilloteau (IRAP), Ting-Yun Cheng (University of Nottingham), Kate Storey-Fisher (NYU), Madeline Lucey (University of Texas Austin), Aurora Cossairt (Argonne National Lab)

%   \item Longer term students: James Butler (University of Chicago), Xiaofeng Dong (University of Chicago), Xin Liu (University of Chicago).

% \end{enumerate}






% \subsection*{Workshop participation}
\section*{Honors, grants \& awards}

\begin{enumerate}


\item Kavli Summer Program fellowship grant, 2019. 

\item Physics and Astronomy Department Scholarship, University of Kansas, 2018. 

\item Division of Astrophysics Travel grant, American Physical Society, 2018. 

\item Graduate Research travel award, University of Kansas, 2017.

\item High Energy Physics - Center for Computational Excellence summer fellowship, 2017.

\item SAMSI Travel grant, Astrophysical population emulation workshop, 2016, 2017.

\item Division of Astrophysics Travel grant, American Physical Society, 2016. 

% \item SAMSI Travel grant, ASTRO workshop, 2016.

\item Graduate Research Competition Award, University of Kansas, 2016.

\item National Science Foundation, the American Physical Society and the Sociedad Mexicana de Fisica travel grant, CAM conference, 2015.

\item Junior research fellowship, Council for Scientific and Industrial Research, Government of India, 2012.

\item Innovation in Science Pursuit for Inspired Research (INSPIRE) fellowship, Department of Science and Technology, Government of India, 2008 - 2012.

\end{enumerate}

\section*{Outreach}

\begin{enumerate}

\item Judge: Illinois Junior Academy of Science Region 6 Fair, Skokie, IL 2023.

\item Member: Rubin Observatory Legacy Survey of Space and Time (LSST)- Dark Energy Science Collaboration (DESC), 2018. 

\item President: Society of Physics Students, The University of Kansas chapter, 2015 - 2018.

\item Founding Member: KUbeSat team - a University of Kansas miniature satellite Mission that won the SPS Chapter Research Award by American Institute of Physics, 2016-2018. 

\item Member: American Physical Society (APS), 2014. 

\item Mentor: Afro-Academic, Cultural, Technological and Scientific Olympics (ACT-SO) initiative by National Association for the Advancement of Colored People (NAACP),  2020-2021. 

\item Chief Organizer: TEDxBITSGoa conference, India, 2011.

\end{enumerate}

\bigskip

% Footer
%\begin{center}
%  \begin{footnotesize}
%    Last updated: \today \\
%    \href{\footerlink}{\texttt{\footerlink}}
%  \end{footnotesize}
%\end{center}

\end{document}